\BookPart{The molecular toolkit}{toolkit}{%
Every molecular computation is built from a handful of physical operations:
strands that pair, populations that react, enzymes that cut, join and copy,
amplification that multiplies a signal, and instruments that turn molecules into
measurements. Part~II teaches each operation as a mechanism with inputs,
outputs, rates and failure modes. It supplies the chemistry and physics that
the rest of the book relies on: strand polarity and complementarity, the
thermodynamics of hybridisation, deterministic and stochastic kinetics, enzyme
specificity, PCR, and the interpretation of gels and fluorescence.

\textbf{Chapters in this part}
\begin{itemize}\setlength{\itemsep}{2pt}
\item \textbf{5\quad DNA chemistry and sequence geometry.} Where does strand direction come from, and how do we compute a reverse complement correctly?
\item \textbf{6\quad Hybridisation thermodynamics.} How much duplex forms at equilibrium, and what sets the melting temperature?
\item \textbf{7\quad Reaction kinetics and stochastic chemistry.} How fast does a population approach equilibrium, and when do individual molecules matter?
\item \textbf{8\quad Enzymes as molecular operators.} What exactly do restriction enzymes, ligases and polymerases change?
\item \textbf{9\quad Amplification and copy-number evidence.} How does PCR copy a chosen interval, and what can a fluorescence curve tell us?
\item \textbf{10\quad Separation, readout and measurement.} What does a gel band or a fluorescence signal actually measure?
\end{itemize}}
